\section{Erd\H{o}s Problem \#333: a zero-density set without an extremely thin additive cover}
\noindent Complete counting construction, 24 September 2026.

\subsection*{1) Statement}
If $A\subseteq\mathbb N$ has asymptotic density zero, must there be $B\subseteq\mathbb N$ with $A\subseteq B+B$ and $B(N)=|B\cap[1,N]|=o(\sqrt N)$? The answer is no. The construction uses only nonnegative summands; allowing arbitrary negative summands would be a different question.

\subsection*{2) A finite set which no small basis covers}
For a large integer $N$, set $k=\lfloor\sqrt N\rfloor$, $m=\lceil N/\log N\rceil$, and $I_N=\{N,\ldots,2N-1\}$. We claim there is an $m$-element $A_N\subseteq I_N$ not contained in $D+D$ for any $D\subseteq\{0,\ldots,2N\}$ with $|D|\le k$.

Choose $A_N$ uniformly among the $m$-element subsets of $I_N$. For fixed $D$, the number of its unordered pairs with repetition is at most $k(k+1)/2$, so for sufficiently large $N$,
\[|(D+D)\cap I_N|\le k(k+1)/2\le3N/5.
\]
The probability that $A_N\subseteq D+D$ is at most $(3/5)^m$: if the intersection has size $u\ge m$, the exact probability is $\binom um/\binom Nm$, and each factor $(u-i)/(N-i)$ is at most $u/N\le3/5$; if $u<m$ the probability is zero.

There are at most $(k+1)(2N+1)^k$ possibilities for $D$. The union bound shows that the probability some such $D$ covers $A_N$ is at most
\[(k+1)(2N+1)^k(3/5)^m.
\]
Its logarithm tends to $-\infty$, because the positive part is $O(\sqrt N\log N)$ whereas the negative part has magnitude at least $\log(5/3)N/\log N$, and $\sqrt N/(\log N)^2\to\infty$. The probability is therefore less than one for every sufficiently large $N$, proving the claim. One may choose the lexicographically first qualifying subset, so no infinite probabilistic assertion is needed.

\subsection*{3) Assemble a zero-density infinite set}
For every sufficiently large integer $j$, apply the lemma with $N_j=4^j$, and put
\[A=\bigcup_j A_{N_j}.
\]
The blocks occupy disjoint intervals $[4^j,2\cdot4^j)$. Their sizes are $m_j=\lceil4^j/(j\log4)\rceil$. We have
\[\sum_{i\le j}m_i=o(4^j).
\tag{1}\]
For example, the terms with $i\le j/2$ contribute $O(4^{j/2}+j)$, while those with $j/2<i\le j$ contribute $O(4^j/j+j)$. If $4^j\le X<4^{j+1}$, then $A(X)\le\sum_{i\le j}m_i=o(4^j)=o(X)$. This proves density zero for all cutoffs, including cutoffs inside a block.

\subsection*{4) Every additive cover violates the desired thinness}
Suppose $A\subseteq B+B$. Every representation of an element of $A_{N_j}\subseteq[0,2N_j]$ uses two nonnegative elements of $B$ at most $2N_j$. Thus $D_j=B\cap[0,2N_j]$ covers $A_{N_j}$ by two-term sums. The finite lemma forces
\[|D_j|>\lfloor\sqrt{N_j}\rfloor.
\]
Consequently
\[\limsup_{X\to\infty}\frac{|B\cap[0,X]|}{\sqrt X}\ge\frac1{\sqrt2}.
\]
Including or excluding zero changes this count by at most one. In particular no covering $B$ has $B(X)=o(\sqrt X)$.

\subsection*{5) Outcome and attribution}
\textbf{SOLVED in the negative.} The finite obstruction, its uniform existence for large scales, density-zero assembly, and obstruction to every proposed thin cover are proved. This is a counting-method treatment of the known negative result; no novelty or independent Lean verification is asserted.

Source: \url{https://www.erdosproblems.com/333}, checked 24 September 2026. The source attributes the negative implication to Theorem 2 of Erd\H{o}s--Newman (1977). That theorem is not taken as an unproved input here; the argument above supplies its own finite obstruction. The special affirmative result for squares is outside the assertion disproved here.

The estimate refers to complete disproof of the stated universal implication.
\subsection*{6) COMPLETION ESTIMATE}
\textbf{COMPLETION: 100\%}
